\documentclass[11pt,a4paper]{article}
\usepackage[margin=26mm]{geometry}
\usepackage[T1]{fontenc}
\usepackage{lmodern,amsmath,amssymb,amsthm,booktabs,xcolor}
\usepackage{xurl,microtype}
\usepackage[colorlinks=true,linkcolor=blue!45!black,urlcolor=blue!45!black,citecolor=blue!45!black]{hyperref}
\newtheorem{theorem}{Theorem}[section]
\newtheorem{lemma}[theorem]{Lemma}
\newtheorem{proposition}[theorem]{Proposition}
\newtheorem{corollary}[theorem]{Corollary}
\newcommand{\pos}[1]{\left(#1\right)_+}
\newcommand{\LP}{L_{\mathrm{raw}}}
\newcommand{\ind}{\mathbf 1}
\newcommand{\A}{\mathcal A}
\newcommand{\Pset}{\mathcal P}
\title{Erd\H{o}s 708: auditing the new dot arguments\\
\large Near-quadratic exclusion, a relative row certificate, and the remaining gap}
\author{Research checkpoint based on the supplied dot sketches}
\date{23 November 2026}
\begin{document}
\maketitle
\begin{abstract}
We complete the supplied two-matching proof excluding positive global
surplus maxima with $M\asymp n^2$, without a restriction on prime support
or weights. We also repair the three/four-occurrence certificate to retain
repeated primes and obtain relative row errors. Under its actual
small-period-core hypothesis, this yields $J_2(X)\le Q(r+\|S\|_\infty)$
and works at the tail scale $n$, including sparse tails. Two supplied
large-prime reductions are verified; a further reduction does not need a
common periodic core. Finally, we give an actual-interval counterexample
to replacing a fixed-bank union discrepancy by the number of roots.
\textbf{The unrestricted estimate $g(n)\le(2+o(1))n$ is not proved.}
The exact $2n$ bound is not claimed either. The missing steps are stated
with their quantifiers; none of the conditional results is used as a
general core normalization.
\end{abstract}

\section{The objective and the inputs carried forward}

Let $I=\{x+1,\ldots,x+M\}$, $x\ge0$, and
\[
 X(a)=\sum_{p\le M}w_pv_p(a),\qquad 0\le w_p\le1,
 \qquad J_c(w)=\sum_{a\le M}\pos{X(a)-c}-\sum_{b\in I}\pos{X(b)-1}.
\]
All target valuations are raw. Fix $c=2+\delta>2$. The exact maximization
identity is
\begin{equation}
 \max_wJ_c(w)=\max_{A\subseteq[2,M]}
       \bigl(\LP(A,I)-c|A|\bigr).                     \label{eq:global}
\end{equation}
At a global maximizer take the entire strict tail
$A=\{a\le M:X(a)>c\}$ and write $n=|A|$. The same $w$ is a raw dual
optimum for this $A$. If $z$ is a paired primal optimum and
$F=\LP(A,I)$, then
\begin{equation}
 F=cn+J_c(w),\quad z_b=1\ (X(b)>1),\quad
 \sum_{b\in I}z_bv_p(b)=D_p(A)\ (w_p>0).              \label{eq:paired}
\end{equation}
Also, for every $U\subseteq A$,
\begin{equation}
 \LP(A\setminus U,I)\le F-c|U|.                      \label{eq:deletion}
\end{equation}
Prices on zero-demand primes may be removed without changing the full
tail or the maximal value.

The preceding checkpoint proved a uniform conditional first moment on
every complete additive prefix tail, using disjoint first-crossing cells
and uniform rough-number estimates. Its proof and the subsequent MGF
derivation are included in the accompanying notes; the rough-number
input is recorded in \cite{Fan}. Positive surplus implies
$\sum_p w_p/(p-1)=O_c(1)$. Consequently there are constants $C_0,B_0$
depending only on $c$ such that, at a positive maximum,
\begin{equation}
 \sum_{a\in A}X(a)\le C_0n,\quad F\le C_0n,\quad
 \sum_{a\in A}\exp\!\left(\frac{X(a)-c}{2B_0}\right)\le2n.
                                                               \label{eq:mgf}
\end{equation}
Jensen's inequality gives the particularly useful consequence
\begin{equation}
 \boxed{\quad
 \sum_{a\in U}X(a)\le ct+2B_0t\log(2n/t),\quad t=|U|>0.
 \quad}                                                       \label{eq:subset}
\end{equation}
Thus \emph{any} $o(n)$ exceptional subset has original weighted demand
$o(n)$, uniformly in $M/n$. The older requirement of an extra
$\log(M/n)$ factor is superseded for these complete tails.

Every prime factor of every source, including a zero-price factor,
satisfies $p<C_0n$. Indeed, deleting a source divisible by
$p\ge F-c+1$ leaves a cover of cost at most $F-c$. Its cofactor bank
has at least $p$ positions and at least one unit of remaining capacity.
Buying that unit supplies $a/p$; a remaining single $p$-unit can be
repaired at cost at most one using full-interval supply. This contradicts
$F-c+2<F$. In particular the number of active rows is
$r\le\pi(C_0n)=O_c(n/\log n)$.

We use Wang's already established ambient gate
$M\ge16n^2\Rightarrow\LP(A,I)\le2n$ and raw-to-integer transfer
\cite{WangMain,WangRaw}. Hence a positive extremum has
\begin{equation}
 cn<M<16n^2.                                           \label{eq:range}
\end{equation}
Rounding is not reproved in this checkpoint.

\section{Completion of the near-quadratic argument}

\begin{theorem}\label{thm:nearquad}
Fix $\delta,\kappa>0$. There is no sequence of positive global maxima
of $J_{2+\delta}$ whose complete-tail sizes $n$ tend to infinity and
whose lengths satisfy $M\ge\kappa n^2$.
\end{theorem}

\subsection{Few hot divisors near the cold threshold}
Put $K=\#\{b\in I:X(b)>1\}\le F\le C_0n$ and $T=M/(K+1)$.
Every integer $d\le T$ is cold, meaning $X(d)\le1$: otherwise its
at least $K+1$ interval multiples would be hot.

\begin{lemma}\label{lem:packing}
If the prefix contains a source of load greater than two, then,
uniformly for $C\ge1$,
\[
 \#\{d\le CT:X(d)>1\}=O(C^7).
\]
\end{lemma}
\begin{proof}
First $K\ge\lfloor M^{1/3}\rfloor$. For a source $a$ with $X(a)>2$,
there is a hot divisor $d\le a^{2/3}$. If a prime $p>a^{2/3}$ divides
$a$, its complementary divisor is hot because $w_p\le1$. Otherwise a
prime or an initial product of prime factors gives a divisor in
$[a^{1/3},a^{2/3}]$; it and its complement both have size at most
$a^{2/3}$, and one is hot. The bank of that divisor proves the claim.

Every hot $d\le CT$ exceeds $T$ and has at least $(K+1)/C-1$ targets.
For two such roots,
\[
 |I\cap d\mathbb Z\cap e\mathbb Z|
 \le\frac{M\gcd(d,e)}{de}+1
 \le\frac{(K+1)\gcd(d,e)}T+1.
\]
Take $L=\lceil4C\rceil$, $\alpha=1/(4L^2)$. If $K$ exceeds a
suitable absolute multiple of $C^2$, then $L$ roots with pairwise
gcd at most $\alpha T$ have, by the first Bonferroni inequality, a
bank union of size at least
\[
 L\frac{K}{2C}-\binom L2(2\alpha K+1)>K,
\]
which is impossible. A maximal separated family therefore has fewer
than $L$ representatives. Every other root $d$ has
$g=\gcd(d,h)>\alpha T$ with a representative $h$. Both $d/g$ and
$h/g$ are at most $C/\alpha=O(C^3)$, and the representative and this
pair of integer quotients determine $d$. This gives $O(C^7)$ roots.
For the remaining range $K=O(C^2)$, the earlier bound on $K$ gives
$M<(K+1)^3$, so $CT<C(K+1)^2=O(C^5)$, proving the same conclusion.
\end{proof}

\subsection{The two source groups}
Fix a large $B$ and then
$B'\ge\max(B,C_0+2,5)$. Define
\[
 \mathcal H=\{h\le B'n:X(h)>1\},\qquad H_0=|\mathcal H|.
\]
When $M\ge\kappa n^2$, Lemma~\ref{lem:packing} gives
$H_0=O_{c,\kappa,B'}(1)$.

Let $A_1$ consist of sources with a divisor in $\mathcal H$.
Discard sources expressible as $a=hk$ with both factors hot and at
most $B'n$. There are at most $H_0^2$ such numbers. This removes all
hot-complement cases: a chosen hot $h>T$ has
$a/h<M/T=K+1\le B'n$. For each remaining source choose a hot carrier
$h\le B'n$ with cold complement. It will earn
\begin{equation}
 X(h)-1\ge X(a)-2.                                    \label{eq:payment}
\end{equation}
For any subset of these sources, its source integers lie in the prefix
union of their chosen carrier banks. Inclusion--exclusion transfers that
union to $I$ with error at most $2^{H_0}$: every multiple count differs
by at most one. Hall's deficiency form therefore matches all but
\begin{equation}
 t_1\le H_0^2+2^{H_0}                                \label{eq:t1}
\end{equation}
sources of the first group to distinct targets.

Let $A_2=A\setminus A_1$ and $G=M/(B'n)$. We have $G\le B'n$.
No divisor of an $A_2$ source lies in $[G,B'n]$, since both that divisor
and its complement would then be cold, contradicting $X(a)>2$.
Every such source has a cold divisor
\[
 \sqrt{Bn}\le u\le B'n.
\]
If there is a prime factor at least $\sqrt{Bn}$, use it and the source
prime bound. Otherwise multiply prime factors, including repetitions,
until the product first reaches $\sqrt{Bn}$; it is less than $Bn$.
The carrier $d=a/u$ is hot and greater than $B'n$, and its bank has
at least $\sqrt{Bn}-1$ positions.

Whenever two current second-group carriers have gcd $g\ge G$, merge
them into $g$. It divides every source assigned to either carrier.
The divisor gap forces $g>B'n$, and the new complement obeys
$a/g<G\le B'n$, so it remains cold. Thus the new carrier is hot,
preserves \eqref{eq:payment}, and enlarges both banks. Repeating leaves
distinct carriers with pairwise gcd less than $G$ and bank sizes
$N_i\ge q_0=\sqrt{Bn}-1$.

\subsection{One second moment and one shared payment budget}
Take a subset of $s$ second-group sources and their $k$ distinct
carriers. Put $V=\sum_iN_i$. Since these distinct sources are prefix
multiples of their assigned carriers, $V\ge s$; also $k\le V/q_0$.
For $i\ne j$,
\[
 |I\cap d_i\mathbb Z\cap d_j\mathbb Z|
 \le\frac GM(N_i+1)(N_j+1)+1.
\]
Consequently the bank multiplicity function $\nu$ satisfies
\[
 \sum_I\nu^2\le V+\alpha_nV^2,
 \qquad
 \alpha_n=\frac1{B'n}(1+1/q_0)^2+\frac1{q_0^2}.
\]
The last term retains all pairwise rounding errors. Cauchy--Schwarz
and the monotonicity of $v/(1+\alpha_nv)$ give
\[
 |\Gamma(S)|\ge\frac{V}{1+\alpha_nV}
 \ge\frac{s}{1+\alpha_ns}\ge s-\alpha_nn^2.
\]
The second-group deficiency is thus at most
\begin{equation}
 t_2\le(1/B'+1/B+o(1))n.                              \label{eq:t2}
\end{equation}

The two matchings may intersect. This does not spend a target twice
without justification. If its carriers are $h\le B'n<d$, their gcd
divides a second-group source and is at most $B'n$. The gap forces
$\gcd(h,d)<G\le T$, hence $X(\gcd(h,d))\le1$. Therefore
\begin{equation}
 X(b)-1\ge X(h)+X(d)-X(\gcd(h,d))-1
 \ge (X(h)-1)+(X(d)-1).                              \label{eq:cross}
\end{equation}
There is at most one assigned source from each group at any target.

Let $t$ be the total number of discarded or unmatched sources.
Equations \eqref{eq:payment} and \eqref{eq:cross} yield the exact bound
\[
 J_c(w)\le\sum_{a\text{ unmatched}}(X(a)-c)-\delta(n-t).
\]
Together with \eqref{eq:subset}, this becomes
\begin{equation}
 \boxed{\quad
 \frac{J_c(w)}n\le
 2B_0\frac tn\log\frac{2n}t-\delta\left(1-\frac tn\right).
 \quad}                                               \label{eq:closing}
\end{equation}
Choose $B,B'$ large enough in terms of $c$, and then $n$ large.
Equations \eqref{eq:t1}--\eqref{eq:t2} make the right side negative.
This proves Theorem~\ref{thm:nearquad}.

\paragraph{Precisely what deteriorates.}
For $\kappa_n=M/n^2$, the same packing proof only gives
\[
 H_0\le C_1(c,B')\kappa_n^{-7},\qquad
 t_1\le H_0^2+2^{H_0}.
\]
It extends to $\kappa_n^{-7}=o(\log n)$, for example
$M\ge n^2/(\log n)^a$ with fixed $a<1/7$. It does not make
$t_1=o(n)$ throughout $M=o(n^2)$. The second-group argument itself
does not deteriorate with $\kappa_n$.

\section{The three/four-occurrence certificate with relative errors}

Here the hypothesis on the core is substantive. It is not inferred from
global optimality.

\begin{theorem}\label{thm:core}
Let $S\ge0$ be an actual periodic function of period $Q$, and set
\[
 X(t)=S(t)+\sum_{p\in\Pset}w_pv_p(t),\quad 0\le w_p\le1.
\]
Assume that every active $p\in\Pset$ exceeds $M^{1/5}$, that
$\gcd(p,Q)=1$, and that $Q=o(\log M)$. Put
$r=\#\{p:w_p>0\}$ and $S_{\max}=\max S$. Then, uniformly in the
shift and the weights, for all sufficiently large $M$ along this regime,
\begin{equation}
 \boxed{J_2(X)\le Q(r+S_{\max}).}                     \label{eq:fullcore}
\end{equation}
The proof retains all source prime powers; it discards no $o(M)$ set.
\end{theorem}

\subsection{Cap comparison and residual payments}
Put $f_p=\min(w_pv_p,1)$, $Y=S+\sum_pf_p$ and $T_0=X-Y$.
The function $T_0$ has a nonnegative prime-power indicator expansion,
so $\sum_{a\le M}T_0(a)\le\sum_IT_0(b)$. Where $T_0>0$, a capped
prime already contributes one; hence
\[
 (X-1)_+=(Y-1)_++T_0,
 \qquad (X-c)_+\le(Y-c)_++T_0.
\]
Thus $J_c(X)\le J_c(Y)$. This does not turn $Y$ into another raw
dual or transfer row balances to it.

Fix a residue modulo $Q$ with $S<1$, and write
\[
 s=1-S,\quad u_p=\min(w_p,s),\quad
 e_{pq}=(u_p+u_q-s)_+,
\]
\[
 L(a)=\Bigl(\sum_p\min(f_p(a),s)-1-s\Bigr)_+,
 \quad R(b)=\Bigl(\sum_p\min(f_p(b),s)-s\Bigr)_+.
\]
In row $p$, allow partners $q\le U_p=\sqrt{M/p}$, and let $g_p(b)$
be the largest $e_{pq}$ among present eligible distinct primes; set it
to zero unless $p\mid b$. Only rows $p\le M^{3/5}$ are needed.
At every target,
\begin{equation}
 \sum_pg_p(b)\le2\Bigl(\sum_{p\mid b}u_p-s\Bigr)_+\le2R(b).
                                                               \label{eq:rowbudget}
\end{equation}
For completeness, remove eligibility restrictions and let $U\ge V$
be the largest two present weights. If $U+V\le s$, all prices vanish.
Otherwise let $\mathcal T$ be the other positive rows, $k=|\mathcal T|$.
The row sum is at most
$U+V-s+\sum_{i\in\mathcal T}u_i+k(U-s)$.
Its difference from $2(U+\sum_{i\in\mathcal T}u_i-s)$ is at least zero,
because
$\sum_{i\in\mathcal T}u_i-V+(k-1)(s-U)\ge0$.

A source has at most four external occurrences. With at most two
distinct external primes its residual is zero. With three distinct
occurrences, every pair price is at least $L$, and each row has an
eligible partner, so its total row payment is at least $3L$.
With four distinct occurrences all pairs are mutually eligible:
$q<M^{2/5}$ whereas the remaining two primes have product greater
than $M^{2/5}$, implying $pq^2\le M$. Every perfect matching of the
four primes has pair prices summing to at least $L$, so the rows pay
at least $2L$.

The previously discarded repeated pattern is $p^2qr$. Its residual is
at most $e_{qr}$ because $\min(2w_p,s)\le1$. The two simple rows
$q,r$ are mutually eligible and pay at least $2L$. Patterns $p^4$,
$p^3q$ and $p^2q^2$ have zero residual. For four-occurrence sources
only simple rows and simple partners will be counted below.

\subsection{A common union for both source types}
Fix a row, a residue, and a price level $z\ge0$. Put $y=M^{1/5}$ and
\[
 \mathcal B=\{q\in\Pset:q\ne p,\ y<q\le U_p,\ e_{pq}>z\},
 \quad h=\sum_{q\in\mathcal B}\frac1q,\quad k=|\mathcal B|,
 \quad N=\frac{M}{Qp}.
\]
If $\mathcal B$ is empty there is nothing to count. Otherwise
$y\le U_p\le M^{2/5}$ and Chebyshev's estimate gives
\begin{equation}
 k\le U_ph,\qquad k\ll U_p/\log M,\qquad
 k^2\ll (M/p)h/\log M.                              \label{eq:kerror}
\end{equation}
For the target union of $pq$ banks in that residue, the first moment is
$Nh+O(k)$ and the second is at most $N(h+h^2)+O(k+k^2)$.
Consequently its cardinality $T_p(z)$ satisfies
\begin{equation}
 T_p(z)\ge(1-o(1))\frac{Nh}{1+h}.                     \label{eq:union}
\end{equation}
This is a relative estimate even for arbitrarily small $h$:
$k/(Nh)\le Q/U_p$ and $k^2/(Nh)\ll Q/\log M$. Also
$h\le\log2+o(1)$.

Let $D_{3,p}(z)$ count the three-distinct-occurrence sources seen in
this row and level. Count all possible partners $q\in\mathcal B$ and
third primes $r>y$, subtracting once when both partners belong to
$\mathcal B$. This counts unordered pairs meeting $\mathcal B$ and
is nonnegative on every integer. Multiple counts and
\eqref{eq:kerror} give
\[
 D_{3,p}(z)\le N\sum_{q\in\mathcal B}\frac1q
                    \sum_{y<r\le M/(pq)}\frac1r
                     -\frac N2h^2+o(Nh).
\]
All rounding errors are $O(QNh/\log M)$, including the subtracted
terms. Put
\[
 x(q)=5\log q/\log M\in[1,2].
\]
Uniform prime reciprocal estimates give the inner sum at most
$\log(4-x(q))+o(1)$.
Since the derivative of $\log(4-x)$ with respect to $\log x$ is
at most $-1/3$, and reciprocal-prime mass in $\log x$ has cumulative
mass at most $t+o(1)$, rearrangement gives
\[
 \sum_{q\in\mathcal B}\frac{\log x(q)}q\ge h^2/2-o(h).
\]
Hence
\begin{equation}
 D_{3,p}(z)\le Nh\bigl(\log3-\tfrac23h+o(1)\bigr).    \label{eq:D3}
\end{equation}
The error remains proportional to $h$, which is essential.

For four-occurrence sources count only simple $p$ and simple partners
$q$. The remaining two primes $r,t>y$ may be equal and obey
$rt\le M/(pq)<y^3$. The unordered reciprocal sum is at most
\[
 \gamma+o(1),\qquad
 \gamma=\frac12\int_1^2\frac{\log(3-x)}x\,dx<0.15.
\]
The diagonal adds only $\frac12\sum_{r>y}r^{-2}=o(1)$. The number
of pair moduli $rt\le V$, with $y^2\le V\le y^3$, is
$O(V/\log M)$, so again all floor errors are $O(QNh/\log M)$.
It follows that
\begin{equation}
 D_{4,p}(z)\le Nh(0.15+o(1)).                          \label{eq:D4}
\end{equation}
The strict integral bound can be checked with a four-panel trapezoid
upper sum for the convex function $\log(3-x)/x$; it is less than
$0.149011$. No numerical optimization is used as a theorem.

Put $A_0=\frac23\log3+0.15$ and $B_0'=4/9$. Direct calculation gives
\[
 \max_{h\ge0}(1+h)(A_0-B_0'h)
 =A_0+\frac{(A_0-B_0')^2}{4B_0'}<0.990304<0.991.
\]
Thus \eqref{eq:union}--\eqref{eq:D4} imply, uniformly in all rows,
residues, levels and weights,
\[
 \tfrac23D_{3,p}(z)+D_{4,p}(z)\le0.995T_p(z)
\]
for sufficiently large $M$ along $Q=o(\log M)$. Integrating over
$z$ and using the source payments and \eqref{eq:rowbudget} proves
\begin{equation}
 \sum_{\substack{a\le M\\a\equiv\rho\ (Q)}}L(a)
 \le0.995\sum_{\substack{b\in I\\b\equiv\rho\ (Q)}}R(b).
                                                               \label{eq:residual}
\end{equation}
The two source types use one common target union and one row budget.

\subsection{Restoring the base and single-prime terms}
Let $B=(S-1)_+$ and
$G_p=(S+f_p-1)_+-B$. For $S\ge1$ put $L=R=0$. Everywhere,
\[
 (Y-1)_+=B+\sum_pG_p+R,\qquad
 (Y-2)_+\le B+\sum_pG_p+L.
\]
The discrepancy of the periodic base is at most $QS_{\max}$.
Within each residue $G_p$ is nondecreasing in $v_p$, starts at zero,
and has total increment mass at most one. Expand it over \emph{all}
prime-power levels. Each corresponding congruence count differs between
the prefix and the interval by at most one, so its total discrepancy is
at most $Q$ per prime, independent of the number of levels. Summing
\eqref{eq:residual} gives the stronger estimate
\[
 J_2(Y)\le Q(r+S_{\max})-0.005\sum_{b\in I}R(b).
\]
The cap comparison proves \eqref{eq:fullcore}.

At a genuine original raw extremum, $r\le\pi(C_0n)$ and
$cn<M<16n^2$. If a justified one-sided cap reduction gives the stated
core with $Q=o(\log M)$ and $S_{\max}\le\log_2M$, then
$Q(r+S_{\max})=o(n)$. But on the original full tail,
\[
 J_2(X)-J_{2+\delta}(X)\ge\delta n,
\]
so positive surplus is impossible in this entire class, even when
$n/M\to0$.

\paragraph{The unproved bridge.}
A nonzero raw valuation function on unbounded integers is not periodic.
A bounded capped small-prime profile can be periodic, but the product
of its moduli need not be $o(\log M)$. No such representation or
small-error reduction has been derived here for arbitrary extremal
weights. The theorem therefore cannot be applied to all optimizers.

\section{Large-prime reductions and what they actually remove}

The accompanying audit gives full details of both supplied reductions.
Their relevant quantitative statements are recorded here so that the
scope of each is explicit.

\subsection{Repeated active cofactors near the tail scale}
Retain weights at least a fixed $b>0$, call their load $Y$, and take
sources $2<Y(a)\le L$. Put $k=\lceil L/b\rceil$,
$P=n/(\log n)^{1/(4k)}$, and $H=(\log n)^{1/(2k)}$.
Choose an active $p\in[P,C_0n]$ and let $d$ be the retained active
part of $a/p$. For each fixed $d$, the number $m_d$ of source copies is
\[
 m_d\le\frac Md\sum_{P\le p\le C_0n}\frac1p
 =\frac Md\,O\!\left(\frac{\log\log n}{\log n}\right).
\]
This includes all inactive cofactors in $a=dpu$.
The original paired rows give $\sum_Iz_bY(b)=O(n)$.
Remove targets with $Y>H$; there are $O(n/H)=o(P)$ of them.
Each remaining bank has $(1-o(1))M/d$ positions, since $M/d\ge P$.
A surviving target has at most $H/b$ retained occurrences and hence
only $O_{b,k}(H^k)=O_{b,k}(\sqrt{\log n})$ possible divisors $d$.
Uniformly spreading all $m_d$ copies over their banks gives point load
\[
 O_{b,k}\!\left(\frac{\log\log n}{\sqrt{\log n}}\right)=o(1).
\]
It proves the claimed matching. More strongly its support has $O(n)$
hot targets, so the number of these sources is itself
\begin{equation}
 O_{b,k}\!\left(n\frac{\log\log n}{\sqrt{\log n}}\right)=o(n).
                                                               \label{eq:nearn}
\end{equation}
Equation \eqref{eq:subset} converts this to $o(n)$ original weighted
demand. Discarded positive weights must still be covered by a valid
normalization budget; the matched divisor is the active part, not
automatically the whole original cofactor.

\subsection{Triple-large sources with a periodic lower load}
Let $\log M/\log P\le K_1=O(1)$ and let $\mathcal R$ consist of
active primes at least $P$ occurring with another such prime in a
source. Deleting the other prime gives a hot $p$-divisible bank with
at least $P$ positions. Exact rows imply
\[
 D_p(A)\ge P,\qquad \sum_{p\in\mathcal R}D_p(A)\le K_1n,
 \qquad |\mathcal R|\le K_1n/P.
\]
Suppose an actual periodic lower load $S$ has period $Q$, the three
active primes $p,q,r\ge P$ are coprime to $Q$, and
$S(a)+w_p+w_q>1$. In a fixed row and residue, the $pq$ banks have
natural lengths $L_q=M/(Qpq)\ge P/Q$ and intersections at most
$L_q/q'+1$. For $\epsilon=nQ/P^2\to0$, their union has size
$(1-O_{K_1}(\epsilon))\sum_qL_q$, whereas all corresponding sources
with a third $r\in\mathcal R$ number at most
$O_{K_1}(\epsilon)\sum_qL_q$. Residues are disjoint and their union
is paid by $D_p(A)$. Summing rows gives
\begin{equation}
 |T|=O_{K_1}(n^2Q/P^2)=o(n).                           \label{eq:tripleperiod}
\end{equation}
The periodic load must be a genuine lower load on the actual targets,
and the third prime must be active. These conditions are needed.

\subsection{A varying-core reduction without a common period}
There is also the following consequence of genuine optimality:
\begin{equation}
 \boxed{\quad
 \#\{a\in A:\sum_{p\ge P,\ w_p>0}v_p(a)\ge3\}
 \ll\frac{n^2M}{P^5}+\frac{nM}{P^4},
 \quad P/n^{3/5}\longrightarrow\infty.
 \quad}                                                       \label{eq:tripleplain}
\end{equation}
Here $M\le16n^2$. For large $n$, $P^4>M$, so each such source has
exactly three large occurrences and a remaining cofactor
$u\le U=M/P^3<P$.

For distinct $p,q,r$, fix $p$ and choose any witness $u_qpqr$ for each
partner $q$. The bank of $pqu_q$ is hot and has natural size
$L_q=M/(pqu_q)\ge P$. Since all $u_q<P$,
\[
 \operatorname{lcm}(pqu_q,pq'u_{q'})
 =pqq'\operatorname{lcm}(u_q,u_{q'}),
\]
so intersections are at most $L_q/q'+1$. There are at most
$K_1n/P$ large active primes, and the union has at least
$(1-O(n/P^2))\sum_qL_q$ positions, bounded by $D_p(A)$.
For each $p,q$ there are at most $|\mathcal R|U$ sources, hence at
most $(|\mathcal R|U/P)L_q$. Summing rows gives the first term of
\eqref{eq:tripleplain}.

For $up^2q$, group by the doubled $p$ and use the same $pqu_q$ banks;
there are only $U$ sources per $p,q$, giving the second term.
For $up^3$, at most $U$ sources per active large prime gives the same
bound. Thus repetitions are included. At $P=n^{3/5}\log n$,
\eqref{eq:tripleplain} is $O(n/(\log n)^5)$ and its original weighted
demand is $O(n\log\log n/(\log n)^5)=o(n)$.
This leaves all lower ranges and does not prove coefficient two.

\section{Exact obstructions to two proposed closing shortcuts}

\subsection{An actual union discrepancy much larger than the root count}
Let $y\to\infty$, $N=y^2$, $\mathcal P_y=\{p\le y\}$ and
$Q_y=\prod_{p\le y}p$. Write
$R(C)=\#\{1\le j\le N:\gcd(C+j,Q_y)=1\}$.
The prefix contains only one and primes above $y$ among these rough
numbers, so the prime number theorem gives
\[
 R(0)=1+\pi(y^2)-\pi(y)
      =(1/2+o(1))N/\log y.
\]
On the other hand, averaging over all $Q_y$ shift classes and applying
Mertens' product theorem gives
\[
 \frac1{Q_y}\sum_{C\bmod Q_y}R(C)
 =N\frac{\varphi(Q_y)}{Q_y}
 =(e^{-\gamma}+o(1))N/\log y.
\]
Both classical asymptotics are stated in \cite{Tao}. Therefore some
actual interval $J$ of length $N$ satisfies
\begin{equation}
 \begin{split}
 &\#\{a\le N:\exists p\le y,\ p\mid a\}
   -\#\{b\in J:\exists p\le y,\ p\mid b\}\\
 &\hspace{12mm}\ge(e^{-\gamma}-1/2+o(1))N/\log y.
 \end{split}                                                \label{eq:unioncounter}
\end{equation}
There are only $\pi(y)\sim y/\log y$ roots. Thus the loss divided
by the root count tends to infinity. Every individual and composite
multiple count is nevertheless exactly $\lfloor N/d\rfloor$ or
$\lfloor N/d\rfloor+1$. Replacing $2^{H_0}$ by $O(H_0)$ without
additional structure is false.

The accompanying note lifts this to a complete tail with hot carriers
and cold complements. Explicitly take
\[
 R_0=\lceil3\log_2\log N\rceil,\quad
 h=2^{\lfloor cR_0\rfloor+1},\quad M=hN,\quad
 w_2=1/R_0,
\]
with other prices zero. The full tail is $A=\{h,2h,\ldots,Nh\}$,
so $n=N$ and $M=o(n^2)$. For sources $hb$ with a factor $p\le y$
and $v_2(b)\le R_0$, choose carrier $hp$. They are hot, their
complements are cold, and all carriers are at most $B'n$ for large
$N$. The fixed-carrier Hall deficit is $\gg N/\log N$, while the
number of all hot integers below $B'n$ is $O(N/(\log N)^3)$.

\emph{This is not a counterexample at a positive global optimum.}
Its hot target count satisfies $K/n\to\infty$, and an adaptive choice
of the common carrier $h$ repairs the matching entirely. It identifies
exactly why fixed carrier selection and full-tail membership alone
cannot supply the needed general estimate.
The displayed deficit is itself $o(n)$: it refutes the stronger
$O(H_0)$ shortcut, not the sufficient general $o(n)$ target.

\subsection{A diffuse core cannot simply be replaced by its mean}
Let $S(t)=\sum_{p\in E}w_pv_p(t)$ be any restriction of genuine
positive-price rows of a paired optimum. Then
$\sum_AS=\sum_Iz_bS(b)$. For any constant $t\ge0$ put
\[
 E_t=\sum_{a\in A}|S(a)-t|+\sum_{b\in I}z_b|S(b)-t|.
\]
The triangle inequality gives the exact lower bound
\begin{equation}
 \boxed{E_t\ge(F-n)t.}                                \label{eq:constantobstruction}
\end{equation}
Indeed, the difference between the two signed centered sums is
$(F-n)t$. Since $F>cn$, an $o(n)$ common selected-budget
approximation requires $t=o(1)$; it then also forces
$\sum_AS=o(n)$. Thus an unspecified $O(n)$ error cannot be absorbed
by the strict $0.005$ margin in the row certificate.

\subsection{A valid local state transfer, with its missing hypothesis}
Let $E$ be primes at most $y$, $\mu_E=\sum_{p\in E}1/p$, and let
$d=\prod_{p\in E}p^{e_p}$ specify an exact state. For $\ell$ coprime
to $E$, put $L=M/(d\ell)$. If $L\ge y$, then
\begin{equation}
 \#\{b\in I:d\ell\mid b,\ v_p(b)=e_p\ (p\in E)\}
 \ge L(1-2\mu_E)-1.                                  \label{eq:state}
\end{equation}
There are at least $L-1$ multiples of $d\ell$; exclude those for which
the quotient has an additional factor from $E$. Their number is at
most $L\mu_E+|E|\le L\mu_E+y\mu_E\le2L\mu_E$.
Different exact states use disjoint target sets. Therefore
$\mu_E=o(1)$ gives a genuine relative transfer without a common
period, whenever the banks have $L\ge y$.

This lemma does not establish $\mu_E=o(1)$ for the entire relevant
core at every optimizer. Existing heavy-state estimates supply it only
below a cutoff depending on the fixed weight threshold. They do not
reach all primes up to $M^{1/5}$ for arbitrary small positive weights,
nor do they justify combining separate state matchings with an
independent external-prime budget.

\section{The remaining general obligation}

The near-quadratic theorem forces any unbounded sequence of positive
global extrema into $M/n^2\to0$. It does not resolve either sparse
$n/M\to0$ or positive-density sequences within that range.
For the two-group route, the outstanding obligation is to pay the
first group $A_1$ with $o(n)$ exceptional source demand, jointly with
the already controlled second group. The payment must use genuine
optimality and adaptive carrier choices, and must control multiway
overlap. Pairwise cold gcd is not enough: the three roots $pq,pr,qr$
with unit prices demand three payments at a target divisible by $pqr$,
whose hinge is only two.

One sufficient, explicitly unproved assertion is a uniform adaptive
carrier assignment for the complete tail: all but $o(n)$ sources have
a carrier $d\mid a$ with $X(a/d)\le1$, assigned to interval multiples
so that at every target the sum of their carrier payments is at most
$X(b)-1$. Such an assertion, combined with \eqref{eq:subset}, would
make \eqref{eq:closing} negative throughout the remaining range.
Neither this assertion nor a direct substitute bounding $J_c$ is
proved here. The periodic-core theorem supplies one valid conditional
payment construction, not the missing general normalization.

A weaker one-sided normalization than common selected-budget $L^1$
approximation would suffice for the core route. Specifically, if one
could construct a load $Y$ covered by Theorem~\ref{thm:core}, with
$Q(r+S_{\max})=o_\delta(n)$, such that
\[
 J_2(X)\le J_2(Y)+\lambda\sum_{b\in I}R_Y(b)+o_\delta(n),
 \qquad \lambda<0.005,
\]
then its residual margin would imply $J_2(X)=o_\delta(n)$, contradicting
$J_2(X)\ge J_{2+\delta}(X)+\delta n>\delta n$.
No such uniform one-sided comparison has been established here.

To obtain the requested unrestricted asymptotic through the surplus
route it would suffice, for each $\delta>0$, to exclude positive
extrema with arbitrarily large full tails. Finite tail sizes are bounded
by \eqref{eq:mgf}; this would give a uniform $J_{2+\delta}\le C_\delta$,
then a uniform raw fractional coefficient two, and finally the integer
asymptotic by \cite{WangRaw}. That last universal exclusion remains
unproved. No exact $2n$ conclusion follows from this checkpoint.

\begin{thebibliography}{9}
\small
\setlength{\itemsep}{1pt}
\bibitem{WangRaw} Ziming Wang, \emph{Interval divisor covers---Erd\H{o}s
Problem 708}, raw transfer theorem, repository source inspected
23 November 2026. Exact source identifiers are in the accompanying manifest.
\url{https://github.com/science-prince/erdos-708-interval-divisor-covers/blob/main/papers/main-paper/raw-transfer.tex}.
\bibitem{WangMain} Ziming Wang, same manuscript, quadratic range
appendix and main paper:
\url{https://github.com/science-prince/erdos-708-interval-divisor-covers/blob/main/papers/main-paper/main.tex}.
\bibitem{Fan} Kai Fan, \emph{Numerically explicit estimates for the
distribution of rough numbers}, Journal of Number Theory 260 (2024),
120--150. \href{https://doi.org/10.1016/j.jnt.2024.01.008}{doi:10.1016/j.jnt.2024.01.008}.
\bibitem{Tao} Terence Tao, \emph{Mertens' theorems}, 11 December 2013
(the two classical inputs to \eqref{eq:unioncounter}).
\url{https://terrytao.wordpress.com/2013/12/11/mertens-theorems/}.
\end{thebibliography}
\end{document}
